\section{Limitations}
\label{sec:limitations}

We group limitations by kind, because they demand different responses: an unrun
experiment is a schedule problem, an unenforced boundary is a security problem, and
a per-family guarantee surface is a design consequence.

\subsection{Limitations of this paper}

\para{The quantitative evaluation has not been run.} This remains the paper's primary
limitation and it is not mitigated by anything else in it. \secref{sec:eval}
specifies the protocol; \tabref{tab:results} contains no measurements. Claims about
control-plane overhead, drain throughput, gate agreement, and rollback latency are
therefore \emph{unsupported} in this draft --- not weakly supported. A reader should
treat the architectural argument as the paper's contribution and the performance
characterization as future work. Eight deterministic fault/interleaving checks are
published separately as structural evidence; they do not change this boundary.

\para{No user study.} \sys's central value claim is that a release becomes
reviewable. Whether operators actually review more effectively with it, or merely
click through a richer interface, is an empirical question about people that this
paper does not address and that no amount of systems measurement answers.

\para{Single-system evidence.} Every architectural claim is grounded in one
implementation. We argue in \secref{sec:arch:claim} that per-family guarantees follow
from the heterogeneity of the artifact types rather than from our particular
engineering, but a single system cannot establish that a different factoring is
impossible.

\subsection{Limitations of the release path}

These three were found by checking the paper's claims against the implementation
rather than by reasoning about the design, and each is a case where an earlier
draft of this paper claimed more than the code delivers.

\para{The prompt fix is not a system-wide release theorem.} Prompt creation and
version authoring are now non-live, and all in-repository prompt-alias call sites
use the intent-first service in \algref{alg:promote}. This establishes a durable
reconciliation obligation before a prompt alias changes. It does not prove that
workflow, skill, knowledge-base, or future family effects have the same ordering;
their guarantees remain those in \tabref{tab:families}. Nor does a repository call-
site inventory prevent an external administrator from changing an \mlflow{} alias
outside CALIBER. The claim is therefore path- and family-scoped, not universal.

\para{A durable intent is not distributed atomicity.} The prompt provider call and
SQL settlement still cannot commit atomically. A crash after alias assignment can
leave an operation in \code{applying}; the difference is that its exact before and
after versions are already durable and the operation is visible through the
release-operations endpoint. Reconciliation observes the provider and settles
\code{applied}, \code{failed}, or \code{reconcile\_required}. This removes silent
divergence but does not remove the need to operate the reconciler.

\para{\mode{Apply} is an operator confirmation, not an independent approval.} The
route requires the \operator{} scope and mints an approval record that is already
in the approved state, with the acting operator recorded as the approver. There is
a sibling \approver{} scope, but nothing requires that the person who authored or
requested a change differ from the person who applies it. This is a deliberate
human decision point placed after an automated gate, and it is real; it is not
separation of duties, and we no longer describe it as one.

\para{Assistant-draft autonomy is a scoped exception, not a replacement for
\mode{Apply}.} The assistant now has structural separation for two draft policies:
\code{agent\_review} records an approver-scoped agent decision by an identity
different from the author, and \code{full\_autonomy} delegates publication to a
third operator-scoped service identity. The decision is bound to the draft
hash/version and rechecked at publication. This does not alter the refinement-job
\mode{Apply} path above, gated goal-plan steps, or the paper's lack of measured
reviewer calibration. The implementation evidence shows fail-closed mechanics, not
that the reviewer agent makes production-quality decisions.

\subsection{Limitations of the architecture}

\para{Guarantees are per-family, and that is load-bearing.} An adopter must read
\tabref{tab:families} rather than a summary sentence. We defend the factoring but
not the cost: the most likely operator error in \sys is assuming a guarantee
transfers across families because the same interface component is mounted in both
places. The shared version panel is the canonical trap.

\para{Two dual-write boundaries cannot be closed.} Object-storage writes and
external registry effects cannot join the caller's SQL transaction
(\figref{fig:dataflow}). Prompt-alias writes are intent-first and idempotently
reconciled, which makes their divergence observable but does not make them atomic.
Other provider and object-storage crossings retain family-specific compensation. A
failed object write can still leave orphaned bytes requiring a sweeper.

\para{In-memory limiters are per process.} The failed-login throttle and the API rate
limiter are per-process token buckets, so their effective ceilings scale with the
worker count (\secref{sec:exec}). An operator who sizes them as though they were
global will provision a limit several times looser than intended. This is a genuine
sharp edge and the documentation-only mitigation is unsatisfying.

\para{Loop capacity is coupled to request capacity.} Because there is no worker tier
(D5), a burst of queued work competes with request serving inside the same process.
A deployment needing to scale drain capacity independently would be better served by
a broker.

\para{A single environment.} There is no
\code{dev}\arrowto\code{staging}\arrowto\code{prod} ladder, so there is no staging
rehearsal for a release. Because the current prompt record is a born-approved
provenance anchor, a future multi-environment mode must rebuild a
pending/approve/reject path with a requester-and-approver distinction rather than
enable an existing one (D9).

\subsection{Limitations of the security model}

These deserve separate treatment because overclaiming here is more harmful than
overclaiming about latency.

\para{Containment is not isolation.} Local stdio MCP execution and the shipped local
tool runner reduce ambient authority --- command and host allowlists, a sanitized
environment, a private working directory, process timeouts, resource limits --- but
retain ambient host filesystem and network authority. They are not a security
boundary against a malicious tool author. Untrusted authorship requires an
operator-supplied OS-enforced backend~\cite{agache2020firecracker,bubblewrap}. Only
an attested managed sidecar or a non-loopback HTTPS endpoint counts as an execution
boundary for production preflight.

\para{Risk tiering is only as good as the classification.} Capabilities declared
\tier{gated} are correctly withheld from \aria{}'s synchronous tool projection,
including the hand-written approval and publication tools. The registry still does
not cover every hand-written operation, so future classification drift remains a
review and contract-test risk (\secref{sec:components}).

\para{Project scoping is not a verified isolation boundary.} Visibility filters are
applied per path. This is not the same as a repository-wide, verified multi-tenancy
guarantee, and \sys should not be deployed as though it were one.

\para{Audit coverage is a property of call sites, not modules.} The transactional
audit guarantee holds where the helper is called on the caller's session. Importing
it into a module does not establish that every mutation in that module is audited,
and we know of no static check that would.

\para{Fail-closed egress does not cover arbitrary user code.} It governs HTTP and MCP
transports. A tool that makes its own network calls from inside the runner is
constrained by the runner's environment, not by the egress allowlist.

\para{The body cap is not flood protection.} Published-service admission counts raw
request chunks against a configured maximum, which is a per-request application
bound. It is not connection-level or IP-level denial-of-service protection, and a
deployment needing that must place it in front.

\para{Provenance anchors are internal records.} They are not cryptographic
attestations in the in-toto or SLSA sense~\cite{torresarias2019intoto,slsa}. They
establish what \sys recorded, not what an independent verifier can prove.

\para{Prompt injection is constrained, not solved.} Bounding what a capability may do
limits the blast radius of injection reached through retrieved or tool-returned
content~\cite{greshake2023injection,perez2022ignore}. It does not prevent the
injection, and \secref{sec:eval} notes that no adversarial evaluation has been run.

\subsection{Limitations of scope}

\sys is not a composition or multi-agent orchestration framework, has no managed
cloud offering, and does not optimize multi-agent bundles. The \aria{} capability
registry is an emerging seam --- \statAriaTools{} hand-written tools plus a
\statAriaProjected{}-capability projection --- rather than achieved parity, and we
count it as future work rather than as a contribution.
